\documentclass[11pt]{article}
\usepackage[utf8]{inputenc}
\usepackage[T1]{fontenc}
\usepackage[turkish]{babel}
\usepackage{amsmath,graphicx,booktabs,hyperref}
% Synthetic paper for samepage's LaTeX tests (fictional study).
\title{Akademik Metinlerde Anlam Kaymasının Ölçülmesi}
\author{Ayşe Yılmaz}
\begin{document}
\maketitle

\begin{abstract}
Makine çevirisi akademik yayıncılıkta giderek daha sık kullanılmaktadır. Bu çalışmada 412 tümce incelenmiş ve uyum katsayısı $\kappa = 0{,}81$ olarak hesaplanmıştır.
\end{abstract}

\section{Giriş}\label{sec:giris}
Sinir ağı tabanlı sistemler daha akıcı çeviriler üretmektedir~\cite{vaswani2017,koehn2010}.
Ne var ki \textbf{akıcılık} doğruluğun güvencesi değildir; \emph{olumsuzluk} ekinin kaybolması
anlamı tümüyle değiştirebilir.\footnote{Örnekler derlemden alınmıştır.} Ayrıntılar
Tablo~\ref{tab:sonuc}'de verilmiştir. Kayıp oranı \%12 düzeyindedir.

% A comment inside the text stays where it is.
Başarım, kesinlik $P$ ve duyarlılık $R$ değerlerinin harmonik ortalaması olan
\begin{equation}
  F_1 = \frac{2PR}{P+R}
  \label{eq:f1}
\end{equation}
ölçütüyle hesaplanmıştır.

\begin{itemize}
  \item Olumsuzluk kaybı en sık görülen hatadır.
  \item Kesinlik kaybı ikinci sıradadır.
\end{itemize}

\begin{table}[ht]
  \centering
  \caption{Anlam kayması türlerinin dağılımı.}
  \label{tab:sonuc}
  \begin{tabular}{lr}
    \toprule
    Tür & Oran (\%) \\
    \midrule
    Olumsuzluk kaybı & 41,2 \\
    Kesinlik kaybı & 28,7 \\
    \bottomrule
  \end{tabular}
\end{table}

\begin{verbatim}
kod: bu metin çevrilmez
\end{verbatim}

\bibliographystyle{apalike}
\bibliography{kaynaklar}
\end{document}
